\documentclass[10pt]{article}

\usepackage[utf8]{inputenc}

\usepackage[T1]{fontenc}

\usepackage{amsmath}

\usepackage{amsfonts}

\usepackage{amssymb}

\usepackage[version=4]{mhchem}

\usepackage{stmaryrd}

\usepackage{bbold}



\begin{document}

среднеквадратичная оштбка $\mathbf{E}|\hat{X}(t)-X(t)|^{2}$ является минимальной в сравнении со всеми другими оценкамі, составленными по значениям рассматриваемого шроцесса в прошлом до момента $s$ (прогнозирование наз. л и н е й н ы м, если ограничиваются линейными оценками).



Одной из первых была поставлена и решена задача линейного прогнозирования стационарной последовательности, имеющая следующий аналог: в пространстве $L_{2}$ интегрируемых в квадрате функций на отрезке $-\pi \leqslant \lambda \leqslant \pi$ найти проекцию функции $\varphi(\lambda) \in L_{2}$ на подпространство, порожденное функциями $e^{i \lambda k} \varphi(\lambda), k=0$, $-1,-2, \ldots$; эта задача получила шиирокое обобщение в теории стационарных случайных прочессов. Примером для приложений может служить задача прогнозирования случайного процесса, возникающего в системе



\[

L X(t)=Y(t), \quad t>t_{0}

\]



с линейным диффференциальным оператором $L$ порядка $l$ и белым пшумом $Y(t), t>t_{0}$ в правой части; здесь наилучший прогноз $\hat{X}(t), t>s$ по значениям в моменты $t_{0} \leqslant t \leqslant s$ при независимых от белого шума начальных значениях $X^{(k)}\left(t_{0}\right), k=0, \ldots, l-1$, может быть дан с помощью решения соответствующего уравнения



\[

L \hat{X}(t)=0, \quad t>s,

\]



с начальными условиями



\[

\widehat{X}^{(k)}(s)=X^{(k)}(s), \quad k=0, l-1 .

\]



Для систем стохастических дифференциальных уравнений задача прогнозирования одних компонент по значениям других наблюдаемых компонент приводит к соответствующим стохастич. уравнениям әкстраполяции.



Лит. см. при ст. Случайных прочессов интерполячия.



СТУЧ Ю. А. Розанов. дача об оценке значения случайного процесса $Z(t)$ в текущий момент $t$ по каким-либо значениям другого, связанного с ним случайного процесса. Напр., речь может идти об оценке стационарного процесса $Z(t)$ по значениям $X(s), s \leqslant t$, стационарно с ним связанного стационарного процесса (см., напр., [1]). Обычно имеют в виду оценку $\hat{Z}(t)$ с наименьшей среднеквадратичной ошибкой $\mathrm{E}|\hat{Z}(t)-Z(t)|^{2}$. Употребление термина «фильтрация» восходит к задаче 0 выделении сигнала из "смеси» сигнала и случайного шума, одна из важных модиф্фикаций к-рой есть задача оптимальной фильтрации в схемө, когда связь $Z(t)$ и $X(t)$ описывается стохастическим дифференциальным уравнением



\[

d X(t)=Z(t) d t+d Y(t), t>t_{0},

\]



где независимый от $Z(t)$ піум представлен стандартным винеровским процессом $Y(t)$.



Широкое распространение в приложениях получил метод фильтрации (м е т о д К а л м а н а - Б ь юс и), применимый к процессам $Z(t)$, к-рые описываются линейными стохастическими дифференциальными уравнениями. Напр., если в указанной выпе схеме



\[

d X(t)=a(t) Z(t) d t+d Y(t)

\]



при нулевых начальных условиях, то



\[

\hat{Z}(t)=\int_{t_{0}}^{t} c(t, s) d X(t)

\]



где весовая функция $c(t, s)$ находится из уравнений:



\[

\begin{gathered}

\frac{d}{d t} c(t, s)=[a(t)-b(t)] c(t, s), t>s, \\

c(s, s)=b(s), \\

\frac{d}{d t} b(t)=2 a(t) b(t)-[b(t)]^{2}+1, t>t_{0}, \quad b\left(t_{0}\right)=0 ;

\end{gathered}

\]



обобщение этого метода на нелинейные уравнения Іриводит к общим стохастич. уравнениям фильтрации (c.1. [2]).



В случае когда



\[

Z(t)=\sum_{k=1}^{n} c_{k} Z_{k}(t)

\]



зависит от неизвестных параметров $c_{1}, \ldots, c_{n}$, интерполяционную оценку $\hat{Z}(t)$ можно дать, оценивая эти параметры по $X(s), \quad s \leqslant t,-$ здесь применим метод наименьших квадратов и его обобщения (см., напр., [3]).



Лит.: [1] Р о з а н о в Ю. А., Стационарные случайные процессы, М., 1963; [2] Л и пи е p P. ШI, III и р я е в А. Н., м о в И. А., Р Р з а н о в Ю. А., Гауссовские случайные процессы, м., 1970. Ю. А. Розанов.



СМЕЖНЫЙ КЛАСС г р у пп п $G$ п о п о д г р у пп е $H$ (л е в ы ї) - множество элементов группы $G$, равное



\[

a H=\{a h \mid h \in H\},

\]



где $a$-нек-рый фиксированный элемент из $G$. С. к. наз. также левосторонним С. к. группы $G$ по подгруппе $I$, определяемым элементом $a$. Всякий левый С. к. определяется любым из своих элементов. $a H=H$ тогда и только тогда, когда $a \in H$. Для любых $a, b \in G$ С.к. $a H$ и $b H$ либо совпадают, либо не пересекаются. Таким образом, группа $G$ распадается на непересекающиеся левые С. к. по подгруппе $H$ - это разложение наз. левосторонни м разложением группы $G$ по под г р у пп е $H$. Аналогично определяются пр а вы е с ме жны е к л а с сы (множества $H a$, $a \in G)$ и п ра вос торонне е ра з ложение г р у п пы $G$ по подг р у пп е $H$. Оба разложения - правостороннее и левостороннее - группы $G$ по $H$ состоят из одного и того же числа классов (в бесконечном случае совпадают мощности множеств этих классов). Это число (мощность) наз. и н д е к с о м п о д г р у п п ы $H$ в группе $G$. Для нормальных делителей левостороннее и правостороннее разложения совпадают, и в этом случае говорят просто о р а з ложен и и группы по е е но р ма льному д е ли и е л ю.



СМЕІІАННАЯ ГРУППА - группа, к-рая содержит как элементы бесконечного порядка, так и отличные от единицы әлементы конечных порядков (см. Порядок әлемента группы). O. А. Иванова.



СМЕIIIАННАЯ ЗАДАЧА - задача для дифференциальных уравнений и систем с частными производными, содержащая начальные условия и краевые условия, а также задача с носителем данных, состоящем как из характеристич., так и нехарактеристических огределенным образом ориентированных многообразий (см. Смешанная $u$ краевая задачи для гиперболических уравнений и систем, Смешанная и краевая задачи для параболических уравнений и систем, Смешанного типа уравнение). С. з. наз. и краевые задачи для эллиптич. уравнений, когда на различных частях границы заданы условия разного рода (см. Краевая задача для әллиптического уравнения). А. М. Нахушев.



СМЕIIIАНАЯ И КРАЕВАЯ ЗАДАЧИ ДЛЯ ГИПЕРБОЛИЧЕСКИХ УРАВНЕНИЙ И СИСТЕМ - Задачи отыскания решений уравнений и систем с частными производными гиперболич. тина, удовлетворяющих на границе области их задания (или ее части) определенным условиям (см. Краевые условия, Начальные условия).



Краевая задача для гиперболич. уравнений и систем, заданных в нек-рой области $D$ евклидова пространства $\mathbb{R}^{n+1}$, наз. с м е ші а н н о й, или н а ч а л ьн о - к р а е в о ї, если искомое решение, наряду с краевыми условиями, должно удовлетворять и начальным или если носитель $\partial D$ граничных данных состоит





\end{document}